% !TeX program = lualatex
\documentclass[11pt,a4paper]{article}

\usepackage{fontspec}
\usepackage[polish]{babel}
\usepackage[a4paper,top=2.5cm,bottom=2.5cm,left=2.5cm,right=2.5cm]{geometry}
\usepackage{microtype}
\usepackage{xcolor}
\usepackage{graphicx}
\usepackage{booktabs}
\usepackage{tabularx}
\usepackage{xltabular}
\usepackage{enumitem}
\usepackage{tcolorbox}
\usepackage{fancyhdr}
\usepackage{titlesec}
\usepackage{float}
\usepackage{hyperref}

% Definicja cudzysłowu dla języka polskiego
\providecommand{\uv}[1]{,,#1''}

% Kolorystyka
\definecolor{primaryBlue}{RGB}{24, 76, 120}
\definecolor{accentRed}{RGB}{180, 40, 40}
\definecolor{accentGreen}{RGB}{34, 139, 34}
\definecolor{boxBg}{RGB}{246, 248, 251}
\definecolor{boxFrame}{RGB}{185, 200, 215}

\hypersetup{
    colorlinks=true,
    linkcolor=primaryBlue,
    urlcolor=primaryBlue,
    citecolor=primaryBlue
}

% Nagłówki i stopki
\setlength{\headheight}{14pt}
\pagestyle{fancy}
\fancyhf{}
\fancyhead[L]{\small\color{gray} System Sygnalizacji ON-AIR -- Instrukcja Użytkownika}
\fancyhead[R]{\small\color{gray} REAPER DAW \& ESP32}
\fancyfoot[C]{\thepage}
\renewcommand{\headrulewidth}{0.4pt}
\renewcommand{\footrulewidth}{0.4pt}

% Formatowanie sekcji
\titleformat{\section}
  {\normalfont\Large\bfseries\color{primaryBlue}}{\thesection}{1em}{}
\titleformat{\subsection}
  {\normalfont\large\bfseries\color{primaryBlue!80!black}}{\thesubsection}{1em}{}

% Inteligentne makro: sprawdza obecność pliku bezpośrednio w podfolderze photos/
\newcommand{\photoPlaceholder}[3]{%
  \begin{figure}[H]
    \centering
    \begin{tcolorbox}[
      colback=boxBg,
      colframe=boxFrame,
      arc=3mm,
      width=0.92\textwidth,
      halign=center,
      boxrule=1.2pt,
      boxsep=4pt,
      left=6pt, right=6pt, top=6pt, bottom=6pt
    ]
      \IfFileExists{photos/#1}{%
        \includegraphics[width=\linewidth,height=0.38\textheight,keepaspectratio]{photos/#1}%
      }{%
        \vspace{0.35cm}
        {\large\bfseries\color{primaryBlue} [ MIEJSCE NA ILUSTRACJĘ / ZRZUT EKRANU ]}\\[0.2cm]
        {\small\bfseries Sugerowany plik:} \texttt{photos/\detokenize{#1}}\\[0.15cm]
        \parbox{0.82\textwidth}{\centering\small \textbf{Wskazówka montażowa:} #2}
        \vspace{0.35cm}
      }%
    \end{tcolorbox}
    \caption{#3}
    \label{fig:\detokenize{#1}}
  \end{figure}
}

\begin{document}

% ==========================================
% STRONA TYTUŁOWA
% ==========================================
\begin{titlepage}
    \centering
    \vspace*{1.2cm}
    
    {\Huge\bfseries\color{primaryBlue} SYSTEM SYGNALIZACJI ON-AIR}\\[0.4cm]
    {\LARGE\bfseries Bezprzewodowy Sterownik Tablicy Studyjnej dla REAPER DAW}\\[1.2cm]
    
    \begin{tcolorbox}[colback=primaryBlue!5,colframe=primaryBlue,arc=4mm,width=0.88\textwidth,halign=center]
        \vspace{0.3cm}
        {\Large\bfseries INSTRUKCJA UŻYTKOWNIKA}\\[0.2cm]
        {\normalsize Montaż, konfiguracja sieciowa i~codzienna obsługa}
        \vspace{0.3cm}
    \end{tcolorbox}
    
    \vspace{1.2cm}
    
    \photoPlaceholder{box_front_cover.jpg}{Wykonaj ostre zdjęcie całego urządzenia z widocznym podwójnym przełącznikiem kołyskowym (czerwony/zielony), diodami LED oraz podłączonym kasetonem ON-AIR.}{Wygląd gotowego urządzenia sygnalizacyjnego ON-AIR.}
    
    \vfill
    
    \begin{tabularx}{0.88\textwidth}{l X}
        \textbf{Wersja dokumentu:} & 1.2 (Wydanie produkcyjne) \\
        \textbf{Platforma sprzętowa:} & ESP32 (Wi-Fi 802.11 b/g/n) \\
        \textbf{Integracja DAW:} & Cockos REAPER v6.x / v7.x (Python ReaScript) \\
        \textbf{Protokół sterowania:} & MQTT v3.1.1 \\
        \textbf{Data opracowania:} & Październik 2026 \\
    \end{tabularx}
    
    \vspace{0.8cm}
\end{titlepage}

\tableofcontents
\newpage

% ==========================================
% 1. WSTĘP I DZIAŁANIE SYSTEMU
% ==========================================
\section{Wstęp i zasada działania systemu}

System sygnalizacji \textbf{ON-AIR} służy do automatycznego oraz ręcznego sterowania podświetleniem tablicy ostrzegawczej w~studiu nagraniowym, lektorskim lub radiowym. Jego zadaniem jest bezprzewodowe informowanie osób przebywających poza kabiną nagraniową o~trwającej rejestracji dźwięku.

\subsection{Kluczowe właściwości systemu}
\begin{itemize}[leftmargin=1.5em]
    \item \textbf{Pełna automatyzacja z~REAPER DAW:} Włączenie nagrywania (\textit{Record}) w~programie REAPER natychmiast załącza przekaźnik i~podświetla tablicę ON-AIR. Zatrzymanie nagrywania wyłącza lampę.
    \item \textbf{Łączność bezprzewodowa Wi-Fi:} Brak kabli sygnałowych pomiędzy reżyserką a~drzwiami studia -- komunikacja oparta jest o~sieć lokalną Wi-Fi i~protokół MQTT.
    \item \textbf{Niezależny tryb manualny (Passthrough):} Sprzętowe wymuszenie załączenia światła niezależnie od komputera i~sieci Wi-Fi (przydatne podczas prób, wywiadów i~prac lektorskich offline).
    \item \textbf{Pamięć stanu (MQTT Retain):} Gdy sterownik zostanie włączony w~trakcie trwającego nagrywania, od razu odczyta aktualny stan z~serwera i~włączy lampę.
    \item \textbf{Portal konfiguracyjny WWW:} Zmiana sieci Wi-Fi oraz adresu serwera odbywa się wygodnie przez telefon lub przeglądarkę internetową.
\end{itemize}

% ==========================================
% 2. ELEMENTY MANIPULACYJNE I WSKAŹNIKI
% ==========================================
\section{Przełączniki i diody sygnalizacyjne}

Na obudowie sterownika umieszczono podwójny przełącznik kołyskowy, diody informacyjne LED oraz ukryty przycisk przywracania ustawień fabrycznych.

\subsection{Podwójny przełącznik kołyskowy}
Urządzenie wyposażone jest w~podwójny przełącznik kołyskowy:
\begin{enumerate}[leftmargin=1.5em]
    \item \textbf{Klawisz czerwony -- Zasilanie główne (\texttt{Power ON/OFF}):}
    \begin{itemize}
        \item \textbf{Pozycja I (ON):} Włącza zasilanie sterownika i~modułu wykonawczego.
        \item \textbf{Pozycja O (OFF):} Całkowite wyłączenie urządzenia.
    \end{itemize}
    
    \item \textbf{Klawisz zielony -- Tryb ręczny (\texttt{Passthrough / Manual Override}):}
    \begin{itemize}
        \item \textbf{Pozycja I (Załączony):} Wymusza stałe świecenie tablicy ON-AIR. W~tym trybie sterownik \textbf{całkowicie wyłącza moduł Wi-Fi} i~nie nasłuchuje poleceń z~komputera. Dioda statusu sieci gaśnie, a~dioda aktywności świeci ciągle. Pozwala to na ręczne włączenie lampy bez uruchamiania komputera.
        \item \textbf{Pozycja O (Wyłączony -- Tryb AUTO):} Urządzenie pracuje w~trybie automatycznym. Łączy się z~Wi-Fi oraz serwerem MQTT, reagując na uruchomienie nagrywania w~programie REAPER.
    \end{itemize}
\end{enumerate}

\photoPlaceholder{switch_panel.jpg}{Zrób zbliżenie na podwójny przełącznik kołyskowy z czytelnym oznaczeniem klawisza czerwonego (Power) oraz zielonego (Passthrough).}{Podwójny przełącznik kołyskowy zasilania i trybu ręcznego.}

\subsection{Diody sygnalizacyjne LED}
Na obudowie znajdują się trzy wskaźniki świetlne:
\begin{itemize}[leftmargin=1.5em]
    \item \textbf{Dioda zasilania (czerwona):} Świeci zawsze wtedy, gdy włączone jest zasilanie sterownika (czerwony klawisz w~pozycji ON).
    \item \textbf{Dioda Statusu Sieci (niebieska):} Informuje o~stanie połączenia z~siecią Wi-Fi oraz serwerem komunikacyjnym MQTT.
    \item \textbf{Dioda Aktywności ON-AIR (zielona):} Świeci zawsze wtedy, gdy załączone jest zasilanie zewnętrznej tablicy świetlnej ON-AIR.
\end{itemize}

\begin{table}[H]
\centering
\renewcommand{\arraystretch}{1.35}
\caption{Zachowanie Diody Statusu Sieci}
\label{tab:status_led}
\begin{tabularx}{\textwidth}{>{\raggedright\arraybackslash\bfseries}p{3.8cm} >{\raggedright\arraybackslash}X >{\raggedright\arraybackslash}p{4.5cm}}
\toprule
\textbf{Sygnał świetlny} & \textbf{Znaczenie / Stan urządzenia} & \textbf{Co należy zrobić?} \\
\midrule
Wyłączona (OFF) & Urządzenie jest wyłączone LUB włączono zielony klawisz trybu ręcznego. & Brak. W~trybie ręcznym Wi-Fi jest celowo wyłączone. \\
\addlinespace[3pt]
Szybkie miganie (5 razy na sekundę) & Trwa łączenie z~siecią Wi-Fi lub serwerem MQTT. & Zaczekaj kilka sekund na nawiązanie połączenia. \\
\addlinespace[3pt]
Powolne miganie (raz na sekundę) & Tryb konfiguracji (uruchomiona własna sieć \texttt{OnAir\_Setup}). & Połącz się ze sterownikiem przez Wi-Fi, aby go skonfigurować. \\
\addlinespace[3pt]
Błyskawiczny stroboskop & Ostrzeżenie! Przycisk Reset jest wciśnięty dłużej niż 3~sekundy. & Puść przycisk, jeśli nie chcesz skasować ustawień Wi-Fi! \\
\addlinespace[3pt]
Światło ciągłe (SOLID ON) & Połączono z~Wi-Fi i~serwerem MQTT. System działa prawidłowo. & Urządzenie gotowe do pracy z~programem REAPER. \\
\bottomrule
\end{tabularx}
\end{table}

\subsection{Przycisk ustawień fabrycznych (Factory Reset)}
Przycisk umieszczono w~zagłębieniu obudowy, aby uniknąć jego przypadkowego wciśnięcia:
\begin{itemize}[leftmargin=1.5em]
    \item \textbf{Krótkie wciśnięcie (poniżej 3~s):} Ignorowane (zabezpieczenie przed przypadkowym dotknięciem).
    \item \textbf{Przytrzymanie przez 3~sekundy:} Dioda statusu zaczyna bardzo szybko mrugać, ostrzegając o~zbliżającym się resecie.
    \item \textbf{Przytrzymanie przez 5~sekund:} Całkowite skasowanie zapisanej sieci Wi-Fi oraz ustawień serwera MQTT. Urządzenie uruchomi się ponownie w~trybie konfiguracji początkowej.
\end{itemize}
\newpage

% ==========================================
% 3. INSTALACJA W REAPER DAW
% ==========================================
\section{Instalacja i konfiguracja w programie REAPER}

Aby program REAPER automatycznie sterował lampą, należy jednorazowo wgrać do niego skrypt integracyjny.

\subsection{Wymagania wstępne}
\begin{enumerate}[leftmargin=1.5em]
    \item Zainstalowany program REAPER (wersja 6.x lub 7.x) oraz środowisko Python w~systemie operacyjnym.
    \item Włączona obsługa skryptów Python w~REAPERze:
    \begin{itemize}
        \item Otwórz w~górnym menu: \texttt{Options $\rightarrow$ Preferences $\rightarrow$ Plug-ins $\rightarrow$ ReaScript}.
        \item Zaznacz opcję \textbf{\uv{Enable Python for ReaScript}}.
        \item Upewnij się, że w~oknie widnieje zielony komunikat potwierdzający wykrycie biblioteki (np. \texttt{python311.dll is installed and working}).
    \end{itemize}
    \item Zainstalowana biblioteka komunikacyjna \texttt{paho-mqtt} (w~wierszu poleceń systemu wpisz: \texttt{pip install paho-mqtt}).
\end{enumerate}

\photoPlaceholder{reaper_python_prefs.jpg}{Zrzut ekranu okna preferencji REAPER: Options -> Preferences -> Plug-ins -> ReaScript z widoczną włączoną obsługą Pythona.}{Weryfikacja obsługi języka Python w programie REAPER.}

\subsection{Kopiowanie plików skryptu krok po kroku}
\begin{enumerate}[leftmargin=1.5em]
    \item W~programie REAPER wybierz z~górnego menu:
    \begin{center}
        \textbf{\texttt{Options $\rightarrow$ Show REAPER resource path in explorer/finder...}}
    \end{center}
    Otworzy się folder z~plikami konfiguracyjnymi programu.
    
    \item Otwórz podfolder o~nazwie \textbf{\texttt{Scripts}}.
    \item Skopiuj plik \textbf{\texttt{onAirSystem.py}} oraz plik autostartu \textbf{\texttt{\_\_startup.lua}} i~wklej je bezpośrednio do folderu \texttt{Scripts}.
    \item \textbf{Pierwsze uruchomienie i~rejestracja skryptu:}
    \begin{itemize}
        \item Wciśnij klawisz \texttt{?} na klawiaturze w~REAPERze (otworzy się okno \textit{Action List}).
        \item Kliknij przycisk \textbf{\texttt{New action... $\rightarrow$ Load ReaScript...}}, wybierz skopiowany plik \texttt{onAirSystem.py} i~kliknij \textbf{Run}.
        \item Dzięki obecności pliku \texttt{\_\_startup.lua} przy każdym kolejnym włączeniu REAPERa system ON-AIR będzie uruchamiał się całkowicie automatycznie w~tle.
    \end{itemize}
\end{enumerate}

\subsection{Konfiguracja połączenia w REAPERze}
Przy pierwszym uruchomieniu skryptu na ekranie pojawi się małe okno konfiguracji:
\begin{itemize}[leftmargin=1.5em]
    \item \textbf{Broker IP:} Adres IP serwera MQTT w~sieci lokalnej (np. \texttt{192.168.0.100}).
    \item \textbf{Port:} Port komunikacyjny (domyślnie: \textbf{\texttt{1883}}).
    \item \textbf{Topic:} Kanał komunikacji (domyślnie: \textbf{\texttt{studio/onair}}).
\end{itemize}

\photoPlaceholder{reaper_dialog_mqtt.jpg}{Zrzut ekranu okna dialogowego "OnAir MQTT Setup" w programie REAPER z wpisanym adresem IP brokera, portem 1883 i tematem studio/onair.}{Okno dialogowe konfiguracji połączenia w REAPERze.}

Po zatwierdzeniu dane zostaną zapamiętane w~pliku \texttt{onAirConfig.json} w~folderze \texttt{Scripts} i~okno nie będzie pojawiać się ponownie. 

\begin{tcolorbox}[colback=boxBg,colframe=boxFrame,arc=2mm]
\small \textbf{Wskazówka:} Jeśli w~przyszłości zmieni się adres IP serwera w~studiu, wystarczy usunąć plik \texttt{onAirConfig.json} z~folderu \texttt{Scripts}. Przy kolejnym uruchomieniu REAPER ponownie zapyta o~nowy adres IP.
\end{tcolorbox}

\newpage

% ==========================================
% 4. KONFIGURACJA SIECIOWA URZĄDZENIA (WIFI)
% ==========================================
\section{Konfiguracja Wi-Fi sterownika}

Nowe lub zresetowane urządzenie nie zna hasła do sieci Wi-Fi w~studiu, dlatego po włączeniu tworzy własną, tymczasową sieć konfiguracyjną.

\subsection{Procedura połączenia krok po kroku}
\begin{enumerate}[leftmargin=1.5em]
    \item \textbf{Włączenie urządzenia:} Upewnij się, że zielony klawisz jest wyłączony (\textbf{O}), a~następnie włącz zasilanie czerwonym klawiszem (\textbf{I}).
    \item \textbf{Sygnalizacja:} Dioda statusu sieci zacznie \textbf{powoli mrugać} (co 1~sekundę). Oznacza to, że sterownik czeka na konfigurację.
    \item \textbf{Połączenie ze sterownikiem:}
    \begin{itemize}
        \item Otwórz ustawienia Wi-Fi w~swoim smartfonie lub laptopie.
        \item Wyszukaj i~połącz się z~otwartą siecią o~nazwie: \textbf{\texttt{OnAir\_Setup}}.
    \end{itemize}
\end{enumerate}

\begin{tcolorbox}[colback=accentGreen!5,colframe=accentGreen,arc=2mm,title=\textbf{Automatyczne wstrzymanie limitu czasu}]
Jeśli po utracie połączenia z~routerem nikt nie połączy się z~siecią \texttt{OnAir\_Setup} przez 2~minuty, sterownik samoczynnie się zrestartuje i~spróbuje ponownie połączyć ze starą siecią studyjną. 
\textbf{W~momencie, gdy połączysz się z~nim telefonem, odliczanie zostaje zatrzymane} -- masz nieograniczony czas na spokojne wpisanie ustawień.
\end{tcolorbox}

\begin{enumerate}[leftmargin=1.5em,resume]
    \item \textbf{Otwarcie panelu w~przeglądarce:}
    \begin{itemize}
        \item Panel powinien otworzyć się automatycznie po połączeniu z~siecią. Jeśli tak się nie stanie, otwórz przeglądarkę internetową i~wpisz w~pasku adresu:
        \begin{center}
            \textbf{\url{http://192.168.4.1}}
        \end{center}
    \end{itemize}
    
    \item \textbf{Wprowadzenie ustawień:}
    \begin{itemize}
        \item Kliknij przycisk \textbf{\uv{Configure WiFi}}.
        \item Z~listy wykrytych sieci wybierz studyjną sieć Wi-Fi i~wpisz do niej hasło.
        \item Poniżej uzupełnij dane serwera (identyczne jak wcześniej w~programie REAPER):
        \begin{itemize}
            \item \textbf{MQTT Broker IP:} Adres IP serwera MQTT.
            \item \textbf{MQTT Broker Port:} \texttt{1883} (domyślnie).
            \item \textbf{MQTT Topic:} \texttt{studio/onair} (domyślnie).
        \end{itemize}
        \item Kliknij przycisk \textbf{\uv{Save}} na dole ekranu.
    \end{itemize}
\end{enumerate}

\photoPlaceholder{wifimanager_portal.jpg}{Zrzut ekranu interfejsu konfiguracyjnego WiFiManager (strona 192.168.4.1) z widocznymi polami sieci Wi-Fi oraz polami MQTT Broker IP, Port i Topic.}{Strona konfiguracyjna sterownika w przeglądarce WWW.}

\subsection{Zakończenie konfiguracji}
Po kliknięciu \uv{Save} sterownik zapisze dane i~uruchomi się ponownie:
\begin{enumerate}[leftmargin=1.5em]
    \item Dioda statusu zamruga szybko (łączenie z~routerem i~serwerem).
    \item Po kilku sekundach dioda zacznie \textbf{świecić światłem ciągłym}.
    \item Urządzenie jest skonfigurowane. Pamięć jest nieulotna -- po wyłączeniu z~prądu sterownik zachowa wszystkie ustawienia.
\end{enumerate}

\newpage

% ==========================================
% 5. CODZIENNA EKSPLOATACJA
% ==========================================
\section{Codzienna obsługa w studiu}

\subsection{Praca z programem REAPER (Tryb automatyczny)}
\begin{enumerate}[leftmargin=1.5em]
    \item Włącz zasilanie sterownika ON-AIR czerwonym klawiszem (zielony klawisz musi być w~pozycji wyłączonej \textbf{O}).
    \item Odczekaj kilka sekund, aż dioda statusu sieci zaświeci światłem ciągłym.
    \item Uruchom program REAPER na komputerze w~reżyserce.
    \item Włączenie nagrywania w~REAPERze (skrót \texttt{Ctrl + R} lub przycisk \textit{Record}) natychmiast zaświeci tablicę ON-AIR przed studiem oraz zieloną diodę na sterowniku.
    \item Zatrzymanie nagrywania (\texttt{Spacja}) lub zamknięcie programu REAPER automatycznie wyłączy tablicę ostrzegawczą.
\end{enumerate}

\subsection{Praca bez komputera / Próba (Tryb ręczny)}
Gdy w~studiu odbywają się próby, wywiady lub nagrania na zewnętrzny rejestrator bez udziału programu REAPER:
\begin{enumerate}[leftmargin=1.5em]
    \item Przestaw zielony klawisz w~pozycję \textbf{I (Załączony)}.
    \item Lampa ON-AIR oraz zielona dioda aktywności zaświecą się na stałe.
    \item Po zakończeniu pracy przestaw zielony klawisz z~powrotem w~pozycję \textbf{O (Wyłączony)}. Sterownik w~ciągu kilku sekund samoczynnie powróci do trybu bezprzewodowego Wi-Fi.
\end{enumerate}

% ==========================================
% 6. ROZWIĄZYWANIE PROBLEMÓW (FAQ)
% ==========================================
\section{Rozwiązywanie problemów (FAQ)}

\renewcommand{\arraystretch}{1.4}
\begin{xltabular}{\textwidth}{>{\raggedright\arraybackslash\bfseries}p{3.8cm} >{\raggedright\arraybackslash}p{4.8cm} >{\raggedright\arraybackslash}X}
\caption{Najczęstsze problemy i sposoby ich rozwiązania} \label{tab:troubleshooting} \\
\toprule
\textbf{Objaw} & \textbf{Prawdopodobna przyczyna} & \textbf{Sposób rozwiązania} \\
\midrule
\endfirsthead
\toprule
\textbf{Objaw} & \textbf{Prawdopodobna przyczyna} & \textbf{Sposób rozwiązania} \\
\midrule
\endhead
\bottomrule
\endfoot
Dioda statusu bez przerwy szybko miga. & Brak zasięgu Wi-Fi lub serwer MQTT w~sieci studyjnej jest wyłączony. & Sprawdź, czy router Wi-Fi działa poprawnie oraz czy serwer z~brokerem MQTT jest uruchomiony. \\
\addlinespace[4pt]
Dioda statusu powoli mruga (raz na sekundę). & Urządzenie przeszło w~tryb konfiguracji (np. po resecie lub utracie sieci). & Połącz się telefonem z~siecią \texttt{OnAir\_Setup} i~skonfiguruj połączenie (Rozdział 4) lub uruchom urządzenie ponownie. \\
\addlinespace[4pt]
Lampa ON-AIR nie reaguje na nagrywanie w~REAPERze (mimo że dioda statusu świeci ciągle). & 1. Niezgodny kanał (\textit{Topic}) pomiędzy REAPERem a~sterownikiem.\newline 2. Skrypt w~REAPERze nie działa w~tle. & 1. Sprawdź w~pliku \texttt{onAirConfig.json}, czy wpisany temat zgadza się z~ustawieniem sterownika (domyślnie \texttt{studio/onair}).\newline 2. Uruchom skrypt ręcznie z~\textit{Action List} (klawisz \texttt{?}). \\
\addlinespace[4pt]
Lampa ON-AIR świeci cały czas, mimo że REAPER nie nagrywa. & Włączony jest zielony klawisz trybu ręcznego (\textit{Passthrough}). & Przestaw zielony przełącznik kołyskowy w~pozycję wyłączoną (\textbf{O}). \\
\addlinespace[4pt]
W~studiu wymieniono router lub zmieniono hasło do Wi-Fi. & Sterownik próbuje łączyć się ze starą, nieistniejącą siecią. & Przytrzymaj ukryty przycisk Reset przez 5~sekund, aż dioda zamruga stroboskopowo i~zgaśnie. Następnie skonfiguruj nową sieć (Rozdział 4). \\
\end{xltabular}

\end{document}